\section{Forensische Video-Datensätze}

Die Evaluation containerbasierter Forensikverfahren und die Gewährleistung der wissenschaftlichen Reproduzierbarkeit erfordern den Rückgriff auf standardisierte Referenzdatensätze. In der Domäne der Video-Quellenerkennung haben sich zwei Dateisammlungen als primäre Benchmarks etabliert, welche unterschiedliche forensische Schwerpunkte abbilden.

Der \textit{VISION}-Datensatz\footnote{\url{https://lesc.dinfo.unifi.it/VISION/}} ist primär für die Quellengeräte- und Markenidentifikation konzipiert \cite{shullani_vision_2017}. Er umfasst über 1.900 Videodateien, die sowohl im originalgetreuen Zustand unmittelbar von verschiedenen Smartphone- und Tablet-Modellen aufgezeichnet als auch über Social-Media-Plattformen wie YouTube und WhatsApp distribuiert wurden \cite[Kap. 1]{shullani_vision_2017}. Bei der Erstellung wurden 29 Smartphonemodelle und 35 individuelle Geräte von 11 unterschiedlichen Herstellern genutzt.\cite[Tab. 5]{shullani_vision_2017} Aufgrund dieser Struktur eignet sich VISION zur Validierung von Signaturen nativer Kameras sowie zur Analyse plattformspezifischer Veränderungen. Die Videodateien sind dabei nach Geräte-ID, Aufnahmeart (\textit{flat}, \textit{indoor}, \textit{outdoor}) sowie Verarbeitungstyp (native Aufnahme oder Distribution über eine der vier Social-Media-Plattformen) strukturiert, wodurch sich Testdateien gezielt nach Gerät, Marke oder Plattform filtern lassen.

Komplementär hierzu fokussiert sich der \textit{EVA-7K}-Datensatz\footnote{\url{https://drive.google.com/drive/folders/1_J04iEbhZTxQgJC7cE3fSG8vkP1by8i1}} auf die Erkennung nachträglicher Integritätsverletzungen und Software-Manipulationen \cite{yang_efficient_2020}. Er enthält rund 7.000 Videos, die systematisch mittels verschiedener Bearbeitungsprogramme (darunter \(ffmpeg\), \(Avidemux\) und \(Adobe Premiere\)) modifiziert wurden \cite[S. 947]{yang_efficient_2020}. EVA-7K erlaubt somit eine Klassifizierung der zur Videoverarbeitung eingesetzten Software-Encoder anhand ihrer verbleibenden Spuren im Metadaten-Container. Die Videos sind dabei primär nach der eingesetzten Bearbeitungssoftware sowie nach dem Vorhandensein einer nachträglichen Manipulation (\textit{tampering}) kategorisiert, was eine gezielte Zuordnung zu den jeweiligen Software-Klassen ermöglicht.

Die etablierten Referenzdatensätze weisen im Kontext moderner forensischer Fragestellungen ein deutliches Alterungsproblem auf. Der VISION-Datensatz wurde bereits im Jahr 2017 veröffentlicht, während der auf ihm aufbauende EVA-7K-Datensatz aus dem Jahr 2020 stammt. Trotz ihres Alters sind diese beiden Datenbanken für die Validierung von strukturellen Container-Analysen derzeit in der wissenschaftlichen Praxis etabliert, da neuere standardisierte Benchmarks fehlen \cite{shullani_vision_2017, yang_efficient_2020}. Da die Erstellung dieser Daten deutlich vor dem technologischen Durchbruch aktueller generativer KI-Modelle erfolgte, enthalten sie keinerlei synthetisch erzeugte MP4-Container und dementsprechend auch keine kryptografischen C2PA-Manifeste, wie sie in \cref{sec:videoquellen} als Identifikationsmerkmal beschrieben wurden. Zur Schließung dieser Lücke wird im Rahmen dieser Arbeit ein eigener, ergänzender Datensatz aus 61 KI-generierten Videos herangezogen, die mit 11 verschiedenen KI-Modellen bzw. Modell-Versionen von 5 Anbietern erstellt wurden und dessen Aufbau und Erhebungsmethodik in \cref{sec:kidatensatz} vorgestellt werden.